\section{Background and a reading guide}
\label{sec:background}

\subsection{Irrationality, transcendence, and approximation are different questions}

An irrationality proof rules out an equality \(\xi=p/q\).
An irrationality-exponent bound rules out an unlimited supply of
exceptionally accurate near-equalities. Transcendence rules out all
polynomial equalities with integer coefficients, but by itself does
not quantify how small a nonzero polynomial value may be. The last
distinction is essential here: the transcendence of \(\pi\) already
implies that \(\pi^k\) is irrational for every positive integer \(k\),
yet it does not supply the numerical exponent in
\eqref{eq:measuremeaning}.

For example, the statement \(\mu(\xi)\le E\) controls infinitely many
denominator scales. A long continued-fraction computation supplies
finitely many very good approximations, and may reveal useful
patterns, but cannot exclude a new pattern at larger denominators.
The present proof therefore constructs auxiliary objects for every
sufficiently large index \(n\), and relates that index to an arbitrary
denominator \(q\). This is the reason for the uniformity requirements
in Sections~\ref{sec:analysis} and~\ref{sec:transfer}.

It is also useful to separate three quantities often called a
``measure.'' A numerical exponent \(E\) specifies the power of \(q\).
A positive prefactor \(c_\varepsilon\) specifies the scale of the lower
bound. A denominator threshold \(q_0(\varepsilon)\) specifies where
that lower bound starts. These notes make \(E\) explicit, but do not
tabulate numerical values for the other two quantities in the nested
construction. Allowing an arbitrary positive exponent slack is part
of the theorem, not a rounding convention.

\subsection{What Pad\'e approximation contributes}

For a power series \(f(x)\), ordinary Pad\'e approximation seeks
polynomials \(P,Q\) for which \(Q(x)f(x)-P(x)\) has a prescribed
high-order zero. Hermite--Pad\'e approximation replaces a single
function by several functions. In our setting these are
\[
1,\log x,(\log x)^2,\ldots,(\log x)^m.
\]
For each row \(h\), the expression
\[
R_{h,n}(x)=A_{h,0,n}(x)+A_{h,1,n}(x)\log x
 +\cdots+A_{h,m,n}(x)(\log x)^m
\]
vanishes to high order at \(x=1\). At a fixed evaluation point this
can make its value small even though its coefficients are large.
The integral representation is an efficient way to impose many
vanishing conditions simultaneously.

The vanishing conditions are linear equations, but solving a large
rational linear system independently at every \(n\) would conceal
the denominators and make nonvanishing difficult to track. The
explicit rational kernel instead exposes its poles, factorials,
and local expansions. These are exactly the objects needed later
for integrality and prime-divisor estimates.

There are \(m+1\) rows because a single small polynomial may vanish
at the algebraic number selected by a rational approximation.
An invertible coefficient matrix ensures that the entire family
cannot vanish there. The exact determinant is therefore an
arithmetic selection tool: it supplies at least one useful row at
every denominator scale.

\subsection{Why evaluate at a complex point?}

The principal logarithm at \(\ii\) is
\(\lambda=\ii\pi/2\), so powers of the logarithm encode powers of
\(\pi\). An integer polynomial evaluated at \(\ii\) has coefficients
in \(\ZZ[\ii]\), a ring whose nonzero elements have modulus at
least one. This discreteness replaces the familiar lower bound
\(|a|\ge1\) for a nonzero ordinary integer.

No irrationality conclusion follows merely from making
\(U_{h,n}(\lambda)\) small: its degree is usually much larger than
the target power \(k\). The next step is to compare it with
\(U_{h,n}(y)\), where
\[
y=\frac{\ii}{2}(p/q)^{1/k}
\]
is algebraic. The comparison is useful because an algebraic value
has a nonzero integral resultant, while the logarithmic value is
small by analysis. The divided difference controls the cost of
moving from \(\lambda\) to \(y\).

For odd powers, an alternative motivation comes from the rational
multiple \(\beta(2r+1)/\pi^{2r+1}\). That identity suggests trying
binary linear forms in \(1,\beta(3)\), but it does not require us
to find them. The present method uses a higher-dimensional system
and a resultant instead. The unsuccessful binary experiments in
Appendix~\ref{app:experiments} explain why this detour was useful.

\subsection{The two parameters that must not be interchanged}

The interval configuration, its multiplicities, and \(m\) are fixed
before the limit \(n\to\infty\). The index \(n\), taken through
multiples of \(L\), determines the number of poles in the rational
kernel. The degree \(m\) determines the number of logarithmic powers
and rows. They have different roles even when both are large.

For instance, in the accepted L99 application \(m=171\), whereas
\(n\) is unbounded. A factor \(171!\) is a constant with respect
to \(n\), so its logarithm divided by \(n\) tends to zero. A
factor \(d_n^{171}\) is not a constant: its exponential rate is
\(171\). Treating these two factors alike can easily reverse the
apparent success of a construction.

When the target power \(k\) changes, we may select another fixed
scale \(s\), hence another fixed \(m\). The all-power theorem means
that this can be done for each \(k\), followed by the separate
limit in \(n\). It does not authorize a simultaneous limit
\(k=k(n)\) in the estimates already proved.

\subsection{A map of the proof}

\begin{figure}[ht]
\centering
\begin{tikzpicture}[
 box/.style={draw,rounded corners=2pt,align=center,text width=64mm,
 minimum height=12mm,inner sep=4pt,font=\small},
 >={Latex[length=2mm]},node distance=7mm and 8mm]
\node[box] (a) {Nested intervals and a rational kernel\\
integral coefficients; exact determinant};
\node[box,right=of a] (b) {Prime valuations and arc cells\\
a common divisor with rate \(\gamma\)};
\node[box,below=of a] (c) {Contour and height estimates\\
smallness \(D\); growth \(G\)};
\node[box,below=of b] (d) {Minimal polynomial and resultant\\
nonzero value at a nearby algebraic point};
\node[box,below=11mm of c,text width=64mm] (e)
 {Choose \(n\) from \(q\)\\
positive gap \(D-(k-1)G\)};
\node[box,right=of e] (f) {Irrationality-exponent bound\\
\(m(D+G)/(D-(k-1)G)\)};
\draw[->] (a)--(b);
\draw[->] (a)--(c);
\draw[->] (b)--(d);
\draw[->] (c)--(d);
\draw[->] (c)--(e);
\draw[->] (d)--(f);
\draw[->] (e)--(f);
\end{tikzpicture}
\caption{The proof separates construction, arithmetic, analysis, and
transfer. A numerical improvement must respect all four parts.}
\label{fig:proofmap}
\end{figure}

A first reading can follow the overview, the L99 ledger in
Section~\ref{subsec:l99}, and the transfer proof. A second reading
can fill in the local coefficient expansion and the contour tails.
The arc partition and the formalization guide are best read after
the purpose of a common divisor is clear. Appendix~\ref{app:reader}
collects notation and short checks that can be carried out by hand.

The main prerequisites are elementary complex analysis, polynomial
algebra over \(\QQ\) and \(\QQ(\ii)\), prime valuations, and
asymptotic estimates. Resultants, confluent Vandermonde matrices,
and interval arithmetic are introduced where they enter the proof.
Knowledge of Lean is not required to read the mathematical argument.
